\scanpage{18}
Je dis que \inserted{le $\Lambda$-module} $V$ se réconstitue \inserted{(au moins si $\Lambda$ est un corps)} à partir de

a) l'ens $P:=P(V)^!$.

b) la famille des \inserted{$\Lambda$-torseurs} $(T_x)_{x\in P\deleted{(V)^!}}$ -- i.e. le \inserted{$\Lambda^*_{P\deleted{\unsure{(V)^!}}}$-}torseur relatif $\coprod_{x\in P(V)^!} T_x=V^*$ sur $P\deleted{(V)}$
% slanted
\margin{équivaut à la donnée de l'ens $V^*$, avec l'opération (libre) de $\Lambda^*$ sur $V^*$ $\big\{$}

c) le $\Lambda^*$-torseur $\omega=\omega_V$

d) la donnée de la relation de disjonction dans $P\deleted{(V)^!}\inserted{=P(V)^!}$ [ou ce qui revient au $=$, dans $V^*$ -- cette relation étant \inserted{bi-}stable par l'opération $\Lambda^*$ (i.e. $u,v$ disjoints $\Rightarrow\lambda u,\mu v$ disjoints pour $\lambda,\mu\in\Lambda^*$)].
% slanted
\margin{Si $\Lambda$ est un corps, cette relation est la relation triviale $x\ne y$, et ne demande pas de donnée suppl. $\big[$}

\deleted{\unsure{d)|e)} le $\Lambda^*$-torseur $\omega_{\unsure{V}}=\omega_V$}

\unsure{e)|f)} la donnée, pour $x,y$ disjoints, de l'isom de $\Lambda^*$-torseurs
\[
T_x\wedge T_y\simeq\omega\qquad\text{noté}\quad
u\mathbin{\underset{\Lambda_{\unsure{*}}}{\wedge}}v\mapsto u\wedge v
\]

En effet, pour \deleted{\unsure{$x,y,z\in$}} $u,v,w$ dans $V^*$, deux à deux disjoints, on trouve
\[
u+v+w\quad\Longleftrightarrow\quad u\wedge v=v\wedge w=w\wedge u
\]
donc on arrive à trouver la loi de composition \inserted{partielle} $(u\unsure{,|+}v)\mapsto u+v$ sur $V^*$ pour $u,v$ \uline{disjoints}.

\deleted{On \unsure{en|a}} Dans le cas où $\Lambda$ est un corps, posant
\[
\left\{\begin{aligned}
V&=V^*\amalg\{0\}\\
D_x&=T_x\cup\{0\}\subset V
\end{aligned}\right.
\]
les $D_x$ sont des droites vectorielles sur $\Lambda$, et on
% archive sheet number: 17
